% E1.4. Conferência numérica: derivations/check/03-desenho-produtos.R (passou em 2026-09-18).
% Tradução para o português em 2026-09-20 (E1.4c, D29): só o idioma da prosa;
% enunciados, hipóteses, provas, constantes e numeração ficam como estavam.
% Notação: docs/notacao.md (congelada em E1.1). Macros: derivations/macros.tex.
% Molde: WALL teórico, Assumption 5.7 (population compatibility) e Lemma S.3
% (restricted eigenvalue, matrix Bernstein) de ms_theo_1.tex / supp_theo_1.tex.
% O que é transposto do WALL: a passagem populacional -> empírica por Bernstein
% matricial (Passo 2b do Lema S.3). O que é novo: a Gram do desenho de produtos
% fatora (ou se minora) através de E[XX' | U], e a condição populacional sobre a
% parte em U é DERIVADA da densidade de U limitada por baixo, não imposta.
% Numeração global dos resultados: pendente do chat principal (TAREFA.md, §5);
% os rótulos internos são os que valem.
\documentclass[11pt]{article}
\usepackage[a4paper,margin=2.6cm]{geometry}
\usepackage[T1]{fontenc}
\usepackage[utf8]{inputenc}
\usepackage[brazil]{babel}
\usepackage{enumitem}
\usepackage{xcolor}
\usepackage[colorlinks=true,linkcolor=blue!60!black,citecolor=blue!60!black]{hyperref}
% Macros do WAFC. Implementa docs/notacao.md (esboço; congela em E1.1).
% Único lugar onde uma macro é definida; os derivations/*.tex fazem % Macros do WAFC. Implementa docs/notacao.md (esboço; congela em E1.1).
% Único lugar onde uma macro é definida; os derivations/*.tex fazem % Macros do WAFC. Implementa docs/notacao.md (esboço; congela em E1.1).
% Único lugar onde uma macro é definida; os derivations/*.tex fazem \input{macros}.
\usepackage{amsmath,amssymb,amsthm}
\newcommand{\R}{\mathbb{R}}
\newcommand{\E}{\mathbb{E}}
\newcommand{\Prob}{\mathbb{P}}
\newcommand{\T}{^{\top}}
\newcommand{\norm}[1]{\left\lVert #1 \right\rVert}
\newcommand{\normn}[1]{\left\lVert #1 \right\rVert_{n}}
\newcommand{\normL}[1]{\left\lVert #1 \right\rVert_{L_2(P)}}
\newcommand{\Besov}[3]{B^{#1}_{#2,#3}}
\newcommand{\bX}{\mathbf{X}}
\newcommand{\bU}{\mathbf{U}}
\newcommand{\bZ}{\mathbf{Z}}
\newcommand{\btheta}{\boldsymbol{\theta}}
\newcommand{\thetastar}{\btheta^{*}}
\newcommand{\bc}{\mathbf{c}}
\newcommand{\PiJ}{\Pi_{J}}
\newcommand{\VJ}{V_{J}}
\newcommand{\NJ}{N_{J}}
\newtheorem{proposition}{Proposition}
\newtheorem{lemma}{Lemma}
\newtheorem{theorem}{Theorem}
\newtheorem{corollary}{Corollary}
\theoremstyle{definition}
\newtheorem{assumption}{Assumption}
\newtheorem{remark}{Remark}
.
\usepackage{amsmath,amssymb,amsthm}
\newcommand{\R}{\mathbb{R}}
\newcommand{\E}{\mathbb{E}}
\newcommand{\Prob}{\mathbb{P}}
\newcommand{\T}{^{\top}}
\newcommand{\norm}[1]{\left\lVert #1 \right\rVert}
\newcommand{\normn}[1]{\left\lVert #1 \right\rVert_{n}}
\newcommand{\normL}[1]{\left\lVert #1 \right\rVert_{L_2(P)}}
\newcommand{\Besov}[3]{B^{#1}_{#2,#3}}
\newcommand{\bX}{\mathbf{X}}
\newcommand{\bU}{\mathbf{U}}
\newcommand{\bZ}{\mathbf{Z}}
\newcommand{\btheta}{\boldsymbol{\theta}}
\newcommand{\thetastar}{\btheta^{*}}
\newcommand{\bc}{\mathbf{c}}
\newcommand{\PiJ}{\Pi_{J}}
\newcommand{\VJ}{V_{J}}
\newcommand{\NJ}{N_{J}}
\newtheorem{proposition}{Proposition}
\newtheorem{lemma}{Lemma}
\newtheorem{theorem}{Theorem}
\newtheorem{corollary}{Corollary}
\theoremstyle{definition}
\newtheorem{assumption}{Assumption}
\newtheorem{remark}{Remark}
.
\usepackage{amsmath,amssymb,amsthm}
\newcommand{\R}{\mathbb{R}}
\newcommand{\E}{\mathbb{E}}
\newcommand{\Prob}{\mathbb{P}}
\newcommand{\T}{^{\top}}
\newcommand{\norm}[1]{\left\lVert #1 \right\rVert}
\newcommand{\normn}[1]{\left\lVert #1 \right\rVert_{n}}
\newcommand{\normL}[1]{\left\lVert #1 \right\rVert_{L_2(P)}}
\newcommand{\Besov}[3]{B^{#1}_{#2,#3}}
\newcommand{\bX}{\mathbf{X}}
\newcommand{\bU}{\mathbf{U}}
\newcommand{\bZ}{\mathbf{Z}}
\newcommand{\btheta}{\boldsymbol{\theta}}
\newcommand{\thetastar}{\btheta^{*}}
\newcommand{\bc}{\mathbf{c}}
\newcommand{\PiJ}{\Pi_{J}}
\newcommand{\VJ}{V_{J}}
\newcommand{\NJ}{N_{J}}
\newtheorem{proposition}{Proposition}
\newtheorem{lemma}{Lemma}
\newtheorem{theorem}{Theorem}
\newtheorem{corollary}{Corollary}
\theoremstyle{definition}
\newtheorem{assumption}{Assumption}
\newtheorem{remark}{Remark}

\newcommand{\bPsi}{\boldsymbol{\Psi}}
\newcommand{\kron}{\otimes}
\newcommand{\lmin}{\lambda_{\min}}
\newcommand{\lmax}{\lambda_{\max}}
\newcommand{\opnorm}[1]{\left\lVert #1 \right\rVert_{\mathrm{op}}}
\DeclareMathOperator{\vecop}{vec}

\title{As matrizes de Gram populacional e empírica do desenho de produtos}
\author{wafc-draft, derivação 03 (E1.4)}
\date{2026-09-18}
\begin{document}
\maketitle

\section*{Objetos}

O modelo e a notação são os de \texttt{docs/notacao.md}. A resposta é
$Y = \sum_{\cv=1}^{p} \fcoef{\cv}(\bU) X_{\cv} + \varepsilon$ com
$\fcoef{\cv}(u) = \lvl{\cv} + \sum_{\md=1}^{q} \gcomp{\cv}{\md}(u_{\md})$, $\bX = (X_1, \ldots, X_p)\T$
e $\bU = (U_1, \ldots, U_q)\T \in [0,1]^q$. Cada $\gcomp{\cv}{\md}$ é expandida nas wavelets
ortonormais periodizadas $\{\wav{\lev}{\tsl} : \lev = 0, \ldots, J-1;\ \tsl = 0, \ldots, 2^{\lev}-1\}$ de uma
família de suporte compacto (D2: $\jzero = 0$, com a função de escala constante descartada), de modo que cada
bloco $(\cv, \md)$ carrega $\NJ = 2^J - 1$ wavelets e a dimensão penalizada é $d = p q \NJ$.

Escreva $\bPsi(u) \in \R^{K}$, $K = 1 + q\NJ$, para o vetor que reúne a constante e as wavelets
das $q$ covariáveis moduladoras,
\[
    \bPsi(u) = \bigl(1,\ \wav{\lev}{\tsl}(u_1)_{\lev\tsl},\ \ldots,\ \wav{\lev}{\tsl}(u_q)_{\lev\tsl}\bigr)\T ,
\]
com as wavelets de cada $u_{\md}$ ordenadas por $\lev$ crescente e, dentro de um nível, por $\tsl$ crescente.
A linha da matriz de desenho é, na ordem de colunas congelada em D12 (primeiro as $p$ colunas não penalizadas
$X_{\cv}$, depois os blocos $(\cv, \md)$ em ordem lexicográfica),
\[
    \bZ = \Pi\,\bigl(\bX \kron \bPsi(\bU)\bigr) \in \R^{p + d}, \qquad p + d = pK,
\]
onde $\bX \kron \bPsi(\bU)$ é o produto de Kronecker, cujas entradas são $X_{\cv}\bPsi_a(\bU)$ com
$\cv$ como índice mais lento, e $\Pi$ é a matriz de permutação fixa que leva para a frente as $p$ entradas
$X_{\cv}\cdot 1$. As matrizes de Gram populacional e empírica são
\[
    \Gram = \E[\bZ\bZ\T], \qquad \Gramhat = \frac1n \sum_{i=1}^{n} \bZ_i\bZ_i\T ,
\]
e, para a parte em $\bU$ sozinha, $\Gram_{\bPsi} = \E[\bPsi(\bU)\bPsi(\bU)\T] \in \R^{K \times K}$.
Como $\Pi$ é uma permutação, $\Gram$ e $\Pi\T\Gram\Pi$ têm o mesmo espectro, e toda afirmação abaixo
sobre autovalores, autovalores restritos ou constantes de compatibilidade é invariante à ordem das
colunas; a ordem só importa para a contabilidade de E2.1.

Para um suporte $S \subseteq \{1, \ldots, p + d\}$ com $|S| = s_0$ e um vetor $\delta \in \R^{p+d}$,
a constante de compatibilidade de uma matriz simétrica $A$ em $S$ é, como em Bühlmann e van de Geer (2011,
Seção~6.2),
\[
    \compat^2(S; A) = \min\Bigl\{ s_0\,\delta\T A\,\delta :\ \norm{\delta_S}_1 = 1,\ \norm{\delta_{S^c}}_1 \le 3 \Bigr\},
\]
e a cota trivial $\delta\T A\delta \ge \lmin(A)\norm{\delta}_2^2 \ge \lmin(A)\norm{\delta_S}_1^2/s_0$ dá
\begin{equation}
    \label{eq:compat-lmin}
    \compat^2(S; A) \;\ge\; \lmin(A) \qquad \text{para todo } S .
\end{equation}
Portanto uma cota inferior para $\lmin$ é uma cota de compatibilidade (e de autovalor restrito) em todo
suporte, com a mesma constante e sem restrição de cone, exatamente como no WALL (o Lema S.3 de lá). É nessa
forma que E1.5 consome a condição.

\section*{Hipóteses}

\begin{assumption}[Covariáveis moduladoras]
\label{ass:U}
$\bU$ toma valores em $[0,1]^q$ e tem densidade conjunta $p_{\bU}$ em relação a Lebesgue com
\[
    0 < c_U \le p_{\bU}(u) \le C_U < \infty \qquad \text{para todo } u \in [0,1]^q .
\]
\end{assumption}

\begin{assumption}[Covariáveis lineares]
\label{ass:X}
$\norm{\bX}_\infty \le B_X$ quase certamente, e existem constantes $0 < \kappa_1 \le \kappa_2 < \infty$ tais que
\[
    \kappa_1 \;\le\; \lmin\bigl(\E[\bX\bX\T \mid \bU]\bigr) \;\le\; \lmax\bigl(\E[\bX\bX\T \mid \bU]\bigr) \;\le\; \kappa_2
    \qquad \text{quase certamente.}
\]
\end{assumption}

\begin{assumption}[Base de wavelets]
\label{ass:psi}
As wavelets $\wav{\lev}{\tsl}$, $\lev \ge 0$, são a periodização a $[0,1]$ de uma família ortonormal de
suporte compacto, de modo que $\{1\} \cup \{\wav{\lev}{\tsl}\}$ é ortonormal em $L_2[0,1]$ sob a medida de
Lebesgue e $\int_0^1 \wav{\lev}{\tsl} = 0$. Existe uma constante $C_\psi$, que depende apenas da família,
tal que
\begin{equation}
    \label{eq:pointwise}
    \sup_{u \in [0,1]} \sum_{\lev=0}^{J-1}\sum_{\tsl=0}^{2^{\lev}-1} \wav{\lev}{\tsl}(u)^2 \;\le\; C_\psi\, 2^{J}
    \qquad \text{para todo } J \ge 1 .
\end{equation}
\end{assumption}

A Hipótese~\ref{ass:U} é D11 lida para a lei conjunta; veja a Observação~\ref{rem:joint}. Na
Hipótese~\ref{ass:X}, a constante $\kappa_1$ é a quantidade de que o passo de identificabilidade
(E1.2) já precisa; quando $\bX$ e $\bU$ são independentes, ela se reduz a $\lmin(\E[\bX\bX\T]) \ge \kappa_1$.
A cota~\eqref{eq:pointwise} vale para toda família de suporte compacto: em cada nível $\lev$, no máximo
$L$ wavelets periodizadas são não nulas num dado $u$, onde $L$ é o comprimento do suporte da wavelet
mãe, e $\norm{\wav{\lev}{\tsl}}_\infty = 2^{\lev/2}\norm{\psi}_\infty$, de modo que o lado esquerdo é no máximo
$L\norm{\psi}_\infty^2 \sum_{\lev < J} 2^{\lev} \le L\norm{\psi}_\infty^2\, 2^{J}$. A conferência numérica encontra
$C_\psi \approx 1.35$ para daublets com filtro de tamanho 8.

\section*{Enunciado}

\begin{proposition}[Matriz de Gram do desenho de produtos]
\label{prop:gram}
Sob as Hipóteses~\ref{ass:U} a~\ref{ass:psi}, com $D = p + d = p(1 + q\NJ)$:
\begin{enumerate}[label=\textup{(\roman*)}, leftmargin=*]
    \item \textup{(A parte em $\bU$.)} $c_U \le \lmin(\Gram_{\bPsi}) \le \lmax(\Gram_{\bPsi}) \le C_U$, para todo $J$.
    \item \textup{(Fatoração sob independência.)} Se $\bX$ e $\bU$ são independentes, então
    \begin{align*}
        \Gram &= \Pi\,\bigl(\E[\bX\bX\T] \kron \Gram_{\bPsi}\bigr)\,\Pi\T , \\
        \lmin(\Gram) &= \lmin(\E[\bX\bX\T])\,\lmin(\Gram_{\bPsi}) \ge \kappa_1 c_U , \\
        \lmax(\Gram) &= \lmax(\E[\bX\bX\T])\,\lmax(\Gram_{\bPsi}) \le \kappa_2 C_U .
    \end{align*}
    \item \textup{(Caso geral.)} Sem independência,
    \[
        \kappa_1\, \lmin(\Gram_{\bPsi}) \;\le\; \lmin(\Gram) \;\le\; \lmax(\Gram) \;\le\; \kappa_2\, \lmax(\Gram_{\bPsi}),
    \]
    portanto $\kappa_1 c_U \le \lmin(\Gram) \le \lmax(\Gram) \le \kappa_2 C_U$.
    Em particular, $\compat^2(S; \Gram) \ge \kappa_1 c_U$ para todo suporte $S$, incluídas as coordenadas
    não penalizadas, por~\eqref{eq:compat-lmin}.
    \item \textup{(Versão empírica.)} Sejam $(\bX_i, \bU_i)$, $i = 1, \ldots, n$, cópias i.i.d.\ de $(\bX, \bU)$ e
    $R_J = p\,B_X^2\,(1 + q\,C_\psi 2^{J})$. Então $\norm{\bZ_i}_2^2 \le R_J$ quase certamente e, para todo $t > 0$,
    \[
        \Prob\bigl(\opnorm{\Gramhat - \Gram} \ge t\bigr)
        \;\le\; 2D \exp\!\left( - \frac{n\,t^2}{2R_J\,(\kappa_2 C_U + 2t/3)} \right) .
    \]
    Em consequência, com $\gamma = \kappa_1 c_U$,
    \[
        \Prob\Bigl( \lmin(\Gramhat) \ge \tfrac{\gamma}{2} \Bigr)
        \;\ge\; 1 - 2D\exp\!\left( - \frac{n\,\gamma^2}{8R_J\,(\kappa_2 C_U + \gamma/3)} \right) ,
    \]
    que tende a $1$ sempre que $p\,q\,2^{J}\log(p\,q\,2^{J})/n \to 0$,
    e nesse evento $\compat^2(S; \Gramhat) \ge \gamma/2$ para todo suporte $S$.
\end{enumerate}
\end{proposition}

\begin{remark}[Alternativa com cone para o passo empírico]
\label{rem:cone}
A parte (iv) transfere o autovalor mínimo cheio, que é o que o regime de sieve
$2^{J_n} \asymp n^{1/(2\seff+1)}$ de E1.6 disponibiliza sem custo: ali $p q 2^{J}\log(pq2^{J})/n \to 0$
assim que $p q$ cresce mais devagar que $n^{2\seff/(2\seff+1)}$, a menos de logaritmos. Se $pq$ fosse maior,
a constante de compatibilidade sozinha ainda pode ser transferida pelo desvio entrada a entrada: por
Bühlmann e van de Geer (2011, Corolário~6.8, Seção~6.12), $\norm{\Gramhat - \Gram}_\infty \le \compat^2(S;\Gram)/(32 s_0)$
implica $\compat^2(S; \Gramhat) \ge \compat^2(S; \Gram)/2$, e a união sobre as $D^2$ entradas, cada uma
limitada por $B_X^2 C_\psi' 2^{J}$ com variância no máximo $B_X^4 C_U C_\psi' 2^{J}$, dá
$\norm{\Gramhat - \Gram}_\infty = O_p\bigl(2^{J/2}\sqrt{\log D/n} + 2^{J}\log D/n\bigr)$; a exigência passa a ser
$s_0^2\, 2^{J}\log(pq2^{J})/n \to 0$, que envolve $pq$ apenas pelo logaritmo. É a forma ``$s_0 \log(pq\NJ)$''
antecipada no plano; não é necessária para as taxas de E1.6, e fica registrada para o caso $pq \gg n^{1/2}$.
\end{remark}

\begin{remark}[Densidade conjunta versus marginal]
\label{rem:joint}
A parte (i) é onde a cota inferior $c_U$ sobre a densidade entra, e é a densidade conjunta em $[0,1]^q$
que se faz necessária. Cotas marginais não bastam: se $U_2 = U_1$ quase certamente, cada marginal é
uniforme, mas $\wav{\lev}{\tsl}(U_1) - \wav{\lev}{\tsl}(U_2) \equiv 0$ e $\Gram_{\bPsi}$ é singular. A cota
inferior conjunta é, assim, a condição de ``não colinearidade entre moduladoras'' de
\texttt{proposta-metodo.md}, §4, item 2, tornada quantitativa: ela exclui que um $U_{\md}$ seja função dos
demais. O WALL impõe diretamente a cota inferior sobre $\lmin$ da Gram populacional (a Hipótese~5.7 de lá)
e deixa a cota inferior da densidade num ``papel apenas motivador''; aqui a mesma cota é derivada da
densidade, com a constante explícita $c_U$.
\end{remark}

\begin{remark}[O que $\kappa_1 > 0$ exclui]
\label{rem:kappa}
A Hipótese~\ref{ass:X} falha exatamente quando alguma combinação linear não trivial de $X_1, \ldots, X_p$ é,
quase certamente, função de $\bU$ num conjunto de probabilidade positiva. O caso canônico é
$X_1 \equiv 1$ e $X_2 = h(U_1)$: então $\E[\bX\bX\T \mid \bU]$ tem posto um, $\kappa_1 = 0$, e a
Gram populacional degenera à medida que $J$ cresce, porque $h(u_1)\wav{\lev}{\tsl}(u_1)$ é aproximada pelo
espaço gerado por $\{1, \wav{\lev'}{\tsl'}(u_1)\}$. A conferência numérica exibe isso com $h(u) = u - 1/2$:
$\lmin(\Gram)$ é $4\times 10^{-4}$ em $J = 2$ e abaixo de $10^{-4}$ em $J \ge 3$. Nada da proposição se perde
ao permitir dependência entre $\bX$ e $\bU$, desde que $\kappa_1 > 0$; a independência compra apenas a
fatoração exata da parte (ii), não uma constante melhor.
\end{remark}

\section*{Prova}

\paragraph{Parte (i).}
Para $a \in \R^{K}$, escreva $h_a(u) = a\T\bPsi(u) = a_0 + \sum_{\md=1}^{q} h_{a,\md}(u_{\md})$ com
$h_{a,\md} = \sum_{\lev\tsl} a_{\md,\lev\tsl}\wav{\lev}{\tsl}$. Sob a Hipótese~\ref{ass:psi}, as funções
$1, \wav{\lev}{\tsl}(u_1), \ldots, \wav{\lev}{\tsl}(u_q)$ são ortonormais em $L_2([0,1]^q, du)$: dentro de uma
coordenada, isso é a ortonormalidade do sistema periodizado; entre coordenadas,
$\int \wav{\lev}{\tsl}(u_{\md})\wav{\lev'}{\tsl'}(u_{\md'})\,du = \int_0^1 \wav{\lev}{\tsl}\int_0^1 \wav{\lev'}{\tsl'} = 0$ para
$\md \ne \md'$, e $\int \wav{\lev}{\tsl}(u_{\md})\,du = 0$, ambas porque toda wavelet tem integral nula.
Portanto $\int_{[0,1]^q} h_a(u)^2\,du = \norm{a}_2^2$ e, pela Hipótese~\ref{ass:U},
\[
    c_U\norm{a}_2^2 \;\le\; \E\bigl[h_a(\bU)^2\bigr] = \int h_a(u)^2\,p_{\bU}(u)\,du \;\le\; C_U\norm{a}_2^2 .
\]
Como $a\T\Gram_{\bPsi}a = \E[h_a(\bU)^2]$, isto é (i).

\paragraph{Parte (iii), primeiro.}
Seja $\delta \in \R^{pK}$, rearranjado como $\Delta \in \R^{p \times K}$, com a linha $\cv$ guardando os $K$
coeficientes ligados a $X_{\cv}$, de modo que $\delta = \vecop(\Delta\T)$ e
$(\bX \kron \bPsi(\bU))\T\delta = \bX\T\Delta\bPsi(\bU)$. Pela propriedade da torre,
\[
    \delta\T\,\Pi\T\Gram\,\Pi\,\delta
    = \E\Bigl[\bigl(\bX\T\Delta\bPsi(\bU)\bigr)^2\Bigr]
    = \E\Bigl[\bPsi(\bU)\T\Delta\T\,\E[\bX\bX\T \mid \bU]\,\Delta\bPsi(\bU)\Bigr] .
\]
A Hipótese~\ref{ass:X} limita a forma quadrática interna entre $\kappa_1\norm{\Delta\bPsi(\bU)}_2^2$ e
$\kappa_2\norm{\Delta\bPsi(\bU)}_2^2$ quase certamente, e
$\E\norm{\Delta\bPsi(\bU)}_2^2 = \sum_{\cv=1}^{p} \Delta_{\cv\cdot}\Gram_{\bPsi}\Delta_{\cv\cdot}\T
= \delta\T(I_p \kron \Gram_{\bPsi})\delta$, cujo valor está entre $\lmin(\Gram_{\bPsi})\norm{\delta}_2^2$ e
$\lmax(\Gram_{\bPsi})\norm{\delta}_2^2$. Portanto
\[
    \kappa_1\lmin(\Gram_{\bPsi})\norm{\delta}_2^2 \;\le\; \delta\T\,\Pi\T\Gram\,\Pi\,\delta \;\le\; \kappa_2\lmax(\Gram_{\bPsi})\norm{\delta}_2^2 ,
\]
e, como $\Pi$ é ortogonal, isto é a primeira linha exibida de (iii); a parte (i) dá a segunda. A
afirmação sobre compatibilidade é~\eqref{eq:compat-lmin}.

\paragraph{Parte (ii).}
Sob independência, $\E[\bX\bX\T \mid \bU] = \E[\bX\bX\T]$, e a mesma conta dá, para todo
$\delta$, $\delta\T\Pi\T\Gram\Pi\delta = \E\bigl[\bPsi(\bU)\T\Delta\T\E[\bX\bX\T]\Delta\bPsi(\bU)\bigr]
= \delta\T(\E[\bX\bX\T] \kron \Gram_{\bPsi})\delta$; duas matrizes simétricas com a mesma forma quadrática são
iguais, o que é a fatoração. Os autovalores de um produto de Kronecker de matrizes simétricas são os
produtos dos autovalores, de modo que $\lmin(\Gram) = \lmin(\E[\bX\bX\T])\lmin(\Gram_{\bPsi})$, e o mesmo para
$\lmax$. Por fim, $\lmin(\E[\bX\bX\T]) \ge \kappa_1$ porque $\lmin$ é côncavo e
$\E[\bX\bX\T] = \E\bigl[\E[\bX\bX\T\mid\bU]\bigr]$ (Jensen); a cota superior é o mesmo argumento com
$\lmax$ convexo. (A independência é usada apenas por meio de $\E[\bX\bX\T \mid \bU] = \E[\bX\bX\T]$; essa
condição mais fraca basta para (ii).)

\paragraph{Parte (iv).}
A norma da linha: $\norm{\bZ_i}_2^2 = \norm{\bX_i}_2^2\,\norm{\bPsi(\bU_i)}_2^2 \le pB_X^2\,(1 + qC_\psi 2^{J}) = R_J$
por~\eqref{eq:pointwise} aplicada a cada uma das $q$ coordenadas. O resto transpõe o Passo~2b do
Lema~S.3 do WALL para parcelas não ponderadas. Ponha $A_i = (\bZ_i\bZ_i\T - \Gram)/n$, independentes, centradas,
simétricas, com $\Gramhat - \Gram = \sum_i A_i$. A amplitude: $\opnorm{\bZ_i\bZ_i\T} = \norm{\bZ_i}_2^2 \le R_J$
e $\opnorm{\Gram} \le \E\opnorm{\bZ\bZ\T} \le R_J$, de modo que $\opnorm{A_i} \le 2R_J/n$. A variância:
$\E[(\bZ\bZ\T - \Gram)^2] = \E[\norm{\bZ}_2^2\,\bZ\bZ\T] - \Gram^2 \preceq R_J\,\Gram$, portanto
$\opnorm{\sum_i \E A_i^2} \le R_J\lmax(\Gram)/n \le R_J\kappa_2 C_U/n$ por (iii). A desigualdade de Bernstein
matricial (Tropp, 2012, Teorema~1.4, aplicada a $\pm\sum_i A_i$) dá
\[
    \Prob\bigl(\opnorm{\Gramhat - \Gram} \ge t\bigr)
    \le 2D\exp\!\left( - \frac{t^2/2}{R_J\kappa_2 C_U/n + (2R_J/n)(t/3)} \right)
    = 2D\exp\!\left( - \frac{n\,t^2}{2R_J(\kappa_2 C_U + 2t/3)} \right) .
\]
Tome $t = \gamma/2$. No evento complementar, a desigualdade de Weyl dá
$\lmin(\Gramhat) \ge \lmin(\Gram) - \gamma/2 \ge \gamma/2$, usando (iii). O expoente é
$-c\,n/R_J$ com $c = \gamma^2/\{8(\kappa_2C_U + \gamma/3)\}$ fixo, e $R_J \asymp pq2^{J}$,
$\log(2D) \asymp \log(pq2^{J})$; a cota de probabilidade tende a zero exatamente quando
$n/(pq2^{J}) - \log(pq2^{J}) \to \infty$, o que é implicado por $pq2^{J}\log(pq2^{J})/n \to 0$ (a
condição mais fraca $n/(pq2^J) \to \infty$ sozinha não controla o prefator $2D$, como o WALL observa). A
afirmação sobre compatibilidade no evento é, de novo,~\eqref{eq:compat-lmin}. \qed

\section*{O que isto não cobre}

\begin{itemize}[leftmargin=*]
    \item \textbf{A fronteira do intervalo.} Com \texttt{boundary = "interval"} (Cohen, Daubechies e Vial,
    1993), a base continua ortonormal em $[0,1]$, de modo que a parte (i) e a prova passam com $\NJ$ e
    $C_\psi$ ajustados, \emph{desde que a constante seja representada uma única vez}. As funções de escala
    CDV no nível $\jzero$ reproduzem constantes, de modo que mantê-las junto com os $X_{\cv}$ não
    penalizados torna $\Gram$ singular; o desenho tem de descartar uma função de escala por bloco, ou impor
    $\int \gcomp{\cv}{\md} = 0$ explicitamente. Essa é a questão de E1.2, não de E1.4; a proposição supõe o
    desenho periódico de D2. Numericamente, o \texttt{WaveBased} instalado exige $\jzero \ge 6$ para a
    fronteira de intervalo com filtro de tamanho 20, de modo que a opção fica fora de alcance em $J = 3$.
    \item \textbf{$\bX$ ilimitado.} A parte (iv) usa $\norm{\bX}_\infty \le B_X$ através da amplitude $R_J$.
    Um $\bX$ sub-gaussiano exigiria um passo de truncamento (ou a desigualdade de Bernstein para matrizes
    sub-exponenciais); as partes (i) a (iii) precisam apenas de segundos momentos.
    \item \textbf{O reescalonamento empírico de $\bU$.} D11 deixa a transformação monótona para
    $[\epsilon, 1-\epsilon]$ fora da teoria. Se o desenho prático usa $\widehat F_{\md}(U_{\md})$ no lugar de
    $U_{\md}$, a Gram muda por uma reparametrização aleatória, e o passo de concentração não se aplica como
    está enunciado.
    \item \textbf{$p$ crescente.} As constantes $\kappa_1, \kappa_2, B_X$ são supostas uniformes em $n$; se
    $p = p_n$, a Hipótese~\ref{ass:X} tem de valer uniformemente e $R_J$ carrega o fator $p$. Nada aqui
    trata $p \gg n$.
    \item \textbf{O aperto da cota em $n$ pequeno.} A cota de (iv) não é informativa em $n = 200$, $D = 30$
    ($R_J\log D/n \approx 1.6$): a conferência encontra $\lmin(\Gramhat)$ em cerca de $0.4$ a $0.5$ de
    $\lmin(\Gram)$, a ordem de Marchenko--Pastur $(1 - \sqrt{D/n})^2 \approx 0.38$, enquanto a constante de
    compatibilidade num suporte de tamanho 7 fica em cerca de $0.65$ a $0.75$ do seu valor populacional. A
    quantidade restrita é a que E1.5 usa, e ela é mais saudável que o autovalor cheio nos tamanhos de
    amostra de E2.
\end{itemize}

\section*{Conferência numérica}

\texttt{derivations/check/03-desenho-produtos.R} (daublets, filtro de tamanho 8; $n = 200$; $J = 3$; $p = 2$ com
$X_1 \equiv 1$; $q = 2$; $\bU$ uniforme em $[0,1]^2$, de modo que $c_U = C_U = 1$ e $\Gram_{\bPsi} = I_{15}$; $D = 30$).
A Gram populacional é calculada por quadratura do ponto médio numa grade $400 \times 400$, com
$\E[\bX\bX\T\mid\bU = u]$ em forma fechada. O cenário A tem $X_2 = 1/2 + \eta$, $\eta \sim \mathrm{Unif}[-1,1]$
independente de $\bU$; o cenário B tem $X_2 = (U_1 - 1/2) + \eta$, de modo que $\kappa_1 = 0.2503$ e
$\kappa_2 = 1.3318$ (extremos dos autovalores pontuais sobre $u_1 \in [0,1]$); o cenário C tem
$X_2 = U_1 - 1/2$, de modo que $\kappa_1 = 0$. O que ele imprimiu em 2026-09-18:
\begin{itemize}[leftmargin=*]
    \item Ortonormalidade de $\{1, \wav{\lev}{\tsl}(u_{\md})\}$ na grade até $3 \times 10^{-5}$ (erro de quadratura);
    $C_\psi(J)$ igual a $1.105$, $1.238$, $1.300$, $1.331$ e $1.347$ para $J = 2, \ldots, 6$.
    \item (ii) Cenário A: $\Gram$ é igual ao produto de Kronecker permutado até $7.6 \times 10^{-13}$;
    $\lmin(\Gram) = 0.2500 = 0.2500 \times 1$, $\lmax(\Gram) = 1.3333 = 1.3333 \times 1$.
    \item (iii) Cenário B: $\lmin(\Gram) = 0.2747 \ge \kappa_1 = 0.2503$; $\lmax(\Gram) = 1.2271 \le \kappa_2 = 1.3318$;
    $\lmin(\E[\bX\bX\T]) = 0.4167 \ge \kappa_1$. A cota inferior está a menos de 10\% do valor verdadeiro.
    \item Cenário C: $\lmin(\Gram) = 4 \times 10^{-4}, < 10^{-4}, < 10^{-4}$ para $J = 2, 3, 4$.
    \item (iv) Em 200 réplicas: posto de coluna cheio em toda réplica em $n = 200$ (o menor
    $\lmin(\Gramhat)$ igual a $0.047$ em A e $0.081$ em B); mediana de $\lmin(\Gramhat)/\lmin(\Gram)$ igual a $0.40$ (A) e
    $0.51$ (B) em $n = 200$, e $0.84$ e $0.96$ em $n = 3200$, onde $\lmin(\Gramhat) \ge \lmin(\Gram)/2$ em todas as
    réplicas; a mediana de $\opnorm{\Gramhat - \Gram}$ cai de $0.83$ para $0.21$ (A) e de $0.67$ para $0.16$ (B), um
    fator $4.1$ e $4.2$ para um $n$ dezesseis vezes maior; $\max_i\norm{\bZ_i}_2^2 \le 67.3 \le R_J = 93.6$.
    \item Compatibilidade num suporte de tamanho 7 (as duas colunas não penalizadas e cinco colunas de wavelet
    sorteadas), calculada exatamente por $2^7$ programas quadráticos (um por padrão de sinais de $\delta_S$):
    população $\compat^2(S;\Gram) = 0.368$ (A) e $0.495$ (B), ambos acima de $\lmin(\Gram)$; em $n = 200$, sobre 20
    réplicas, mediana $0.240$ (A) e $0.363$ (B), mínimo $0.182$ e $0.282$, sempre acima de $\lmin(\Gramhat)$.
\end{itemize}

\section*{Referências}

\begin{itemize}[leftmargin=*]
    \item Bühlmann, P.\ e van de Geer, S. (2011). \emph{Statistics for High-Dimensional Data}. Springer.
    (\texttt{buhlmann2011statistics}, verificada em L1.) A numeração foi conferida em L3: o Corolário~6.8 está
    na p.~152, na Seção~6.12, e afirma que $32\,s\,\lVert\Sigma_1-\Sigma_0\rVert_\infty/\compat^2(S;\Sigma_0) \le 1$
    transfere a condição de compatibilidade de $\Sigma_0$ para $\Sigma_1$.
    \item Cohen, A., Daubechies, I.\ e Vial, P. (1993). Wavelets on the interval and fast wavelet transforms.
    \emph{Applied and Computational Harmonic Analysis} 1, 54--81. (\texttt{cohen1993wavelets}, verificada em L1.)
    \item Tropp, J.~A. (2012). User-friendly tail bounds for sums of random matrices. \emph{Foundations of
    Computational Mathematics} 12, 389--434. (\texttt{tropp2012user}, copiada da bibliografia do WALL e
    verificada em L3.)
\end{itemize}

\end{document}
